%=====================================================================
\chapter{The Stack: Setting It Up and Keeping It}\label{app:stack}
%=====================================================================
\normalfigures

Everything in this book from \chref{stack} onward assumes a shop you can reach
in a browser. This appendix gets you one, in a single sitting, and lists what
to do when it refuses --- which on a university laptop it sometimes will.

\section{What You Install}

\textbf{One thing: a container runtime.} Not a web server, not PHP, not a
database. Those three run inside containers and never touch your own machine,
which is why the shop is identical on every laptop in the class and why
uninstalling it is one command.

\begin{tight}
  \item \textbf{Windows or macOS:} Docker Desktop, from
        \texttt{docker.com}. On Windows it needs WSL~2, which its installer
        offers to set up.
  \item \textbf{Linux:} the \texttt{docker.io} and \texttt{docker-compose-v2}
        packages from your distribution, then add yourself to the
        \texttt{docker} group and log out and in again.
\end{tight}

Check it with:

\begin{lstlisting}[language=bash]
docker --version
docker compose version
docker run --rm hello-world
\end{lstlisting}

The third is the one that matters. The first two only prove the program is
installed; the third proves the \emph{engine} is running, which is a different
claim and the one that usually fails.

\section{Starting the Shop}

\begin{lstlisting}[language=bash]
cd code
cp .env.example .env
docker compose up -d --build
python waitfor.py
\end{lstlisting}

The first run takes a few minutes and downloads about 1.5~GB; later runs take
seconds. \chref{stack} measures both and explains why the difference matters.

\rowcolors{2}{white}{lightbg}
% The middle column must hold "http://localhost:8090/admin/" unbroken at
% 10.95 pt typewriter, which needs 5.8 cm; at 5.4 it overflows by 7.3 pt.
\begin{longtable}{@{}L{2.9cm}L{5.9cm}L{6.0cm}@{}}
\toprule
\rowcolor{primary!15}
\textbf{What} & \textbf{Where} & \textbf{Sign in with} \\
\midrule
\endhead
The storefront & \texttt{http://localhost:8090/} & nobody; you are a customer \\
The admin panel & \texttt{http://localhost:8090/admin/} & \texttt{admin} /
\texttt{admin123} \\
The database console & \texttt{http://localhost:8091/} & server \texttt{db},
user \texttt{shop}, password \texttt{shoppw}, database \texttt{opencart} \\
\bottomrule
\end{longtable}

\begin{alertbox}[Those credentials are deliberately bad]
\texttt{admin} / \texttt{admin123} is exactly the kind of credential
\chref{security} goes looking for and finds. It is the default so that you can
measure the shop before and after hardening it; Chapter~16 asks you to change
it, and until then every student in the class can sign in to every other
student's shop if they can reach it.

\smallskip
Which they cannot, because nothing here is published to the internet. If you
ever expose this stack beyond your own machine, change the credentials in
\texttt{.env} \emph{first}.
\end{alertbox}

\section{The Four Commands Worth Memorising}

\rowcolors{2}{white}{lightbg}
\begin{longtable}{@{}L{5.0cm}L{10.0cm}@{}}
\toprule
\rowcolor{primary!15}
\textbf{Command} & \textbf{What it does} \\
\midrule
\endhead
\texttt{docker compose up -d} & Start everything. Safe to run when it is
already running \\
\texttt{docker compose down} & Stop everything. \textbf{Your shop is kept} \\
\texttt{docker compose logs -f shop} & Watch the shop's log. This is where PHP
errors go, and they go nowhere else \\
\texttt{docker compose down -v} & Stop everything \textbf{and delete your
shop}. The reset button, and the one mistake with no undo \\
\bottomrule
\end{longtable}

\section{When It Will Not Start}

Five failures account for nearly all of them, and four are not really about
Docker.

\begin{enumerate}[leftmargin=2.2em, itemsep=5pt, topsep=4pt]
  \item \textbf{``Cannot connect to the Docker daemon''} --- the engine is not
        running. On Windows and macOS, start Docker Desktop and wait for its
        icon to stop animating; the command line is ready several seconds
        before the whale stops moving.
  \item \textbf{``Ports are not available''} or ``address already in use'' ---
        something else on your machine already holds port 8090 or 8091. This is
        common; the development machine for this book had both 8080 and 8081
        taken, which is why the defaults are 8090 and 8091 in the first place.
        Edit \texttt{.env}, change \texttt{SHOP\_PORT}, and bring it up again.
        Change only the port in \texttt{.env} --- \texttt{OC\_HTTP\_SERVER}
        reads the same variable, so links keep working.
  \item \textbf{The shop answers but every page is an error} --- the install
        did not finish. Read \texttt{docker compose logs shop} and look for
        ``SUCCESS! OpenCart successfully installed''. If it is absent, the
        database was not ready; \texttt{docker compose down -v} and start
        again.
  \item \textbf{WSL 2 is not installed, or virtualisation is disabled} ---
        Windows only. Docker Desktop's installer handles the first; the second
        is a BIOS setting (often called ``Intel VT-x'', ``AMD-V'' or
        ``SVM Mode'') that some laptops ship disabled.
  \item \textbf{It worked yesterday and today it does not} --- almost always
        disk space. \texttt{docker system df} will tell you;
        \texttt{docker system prune} reclaims what is not in use. It does not
        touch named volumes, so your shop survives --- but read what it
        proposes before agreeing.
\end{enumerate}

\begin{conceptbox}[Read the log before changing anything]
The single most useful habit in this appendix. A shop that is broken is
broken for a reason, and the reason is almost always written down in
\texttt{docker compose logs shop}.

\smallskip
Restarting, rebuilding and resetting are the three things students try first,
and all three destroy the evidence. Look at the log, then act.
\end{conceptbox}

\section{Backing Up, and Why You Will Wish You Had}

Your shop lives in two named volumes. \texttt{down -v} deletes them, and so
does an afternoon of clearing disk space without reading the prompt.

\begin{lstlisting}[language=bash]
# Dump the database into the repository, where git can see it.
docker compose exec -T db mariadb-dump -ushop -pshoppw opencart \
  > backups/shop-$(date +%Y%m%d).sql

# Put it back.
docker compose exec -T db mariadb -ushop -pshoppw opencart \
  < backups/shop-20261001.sql
\end{lstlisting}

Take one after every chapter's laboratory. It costs two seconds, and
\chref{deploy} measures what the alternative costs.

\begin{alertbox}[Where to keep your own work]
Anything you write that must survive belongs in \texttt{code/}, in the
repository --- not inside a running container, where it is deleted at the next
rebuild, and not only in the admin panel, where it is deleted by
\texttt{down -v}.

\smallskip
This matters most in \chref{modules}, where you write a module. Write it in
\texttt{code/modules/} and mount it; do not edit it in place inside the
container. Students who learn this in Week~14 learn it the expensive way.
\end{alertbox}

\section{The Machine Every Measurement Was Taken On}

Every figure in a \faIcon{tachometer-alt}~\textbf{Measured} box in this book was
produced on this machine. Yours will differ; the shapes and ratios should not.

\rowcolors{2}{white}{lightbg}
\begin{longtable}{@{}L{4.6cm}L{10.4cm}@{}}
\toprule
\rowcolor{primary!15}
\textbf{Item} & \textbf{Value} \\
\midrule
\endhead
Host & Windows 11 Pro, 4 CPU cores, 7.7 GiB available to the container
runtime \\
Container runtime & Docker 28.1.1, Linux containers \\
Web and application & Apache 2.4.68, PHP 8.2.34 with OPcache enabled \\
Database & MariaDB 11.8.9 \\
Application & OpenCart 4.1.0.4, installed from the pinned release in
\texttt{code/shop/Dockerfile} \\
Everything on & one machine, so no figure here includes real network latency \\
\bottomrule
\end{longtable}

\begin{alertbox}[The scope of Chapter 16's security work]
\chref{security} asks you to attack a shop: SQL injection, cross-site
scripting, missing security headers, default credentials, an exposed
administrative path. Every one of those exercises is scoped, in the chapter and
here, to \textbf{the container running on your own machine, reached at
\texttt{localhost}, which you installed and which contains no real person's
data}.

\smallskip
That is the whole of it. Running the same checks against any system you do not
own and have not been authorised in writing to test is a criminal offence in
Pakistan under the Prevention of Electronic Crimes Act 2016, and a sackable one
nearly everywhere else. The skill is valuable precisely because it is
dangerous, and the boundary is not a formality.

\smallskip
\chref{legal} returns to this with the statute in front of it.
\end{alertbox}
